\documentclass[12pt,a4paper]{article}
\usepackage[UTF8,fontset=windows]{ctex}
\usepackage[top=2.5cm,bottom=2.5cm,left=3.0cm,right=2.5cm,headheight=15pt]{geometry}
\usepackage{amsmath,amssymb,amsthm,bm,mathtools}
\usepackage{graphicx}
\usepackage{tikz}
\usetikzlibrary{shapes.geometric,arrows.meta,positioning,fit,calc,backgrounds,decorations.pathreplacing}
\usepackage{xcolor}
\definecolor{natblue}{RGB}{0,63,127}
\definecolor{natblue2}{RGB}{0,100,180}
\definecolor{natgray}{RGB}{90,90,90}
\definecolor{natgreen}{RGB}{0,120,60}
\definecolor{natorange}{RGB}{200,80,0}
\definecolor{natred}{RGB}{160,0,0}
\definecolor{lightblue}{RGB}{220,235,250}
\definecolor{lightgreen}{RGB}{215,240,215}
\definecolor{lightyellow}{RGB}{255,250,210}
\definecolor{lightorange}{RGB}{255,230,200}
\definecolor{lightgray}{RGB}{245,245,245}
\usepackage{titlesec}
\titleformat{\section}{\Large\bfseries\color{natblue}}{\thesection}{0.8em}{}[{\color{natblue}\titlerule[0.8pt]}]
\titleformat{\subsection}{\large\bfseries\color{natblue!85!black}}{\thesubsection}{0.7em}{}
\titleformat{\subsubsection}{\normalsize\bfseries\color{natblue!65!black}}{\thesubsubsection}{0.6em}{}
\titlespacing*{\section}{0pt}{20pt}{10pt}
\titlespacing*{\subsection}{0pt}{14pt}{6pt}
\titlespacing*{\subsubsection}{0pt}{10pt}{4pt}
\usepackage{booktabs,array,tabularx,multirow,longtable,colortbl}
\newcolumntype{C}[1]{>{\centering\arraybackslash}p{#1}}
\newcolumntype{L}[1]{>{\raggedright\arraybackslash}p{#1}}
\newcolumntype{R}[1]{>{\raggedleft\arraybackslash}p{#1}}
\usepackage{enumitem}
\setlist[itemize]{leftmargin=1.8em,itemsep=2pt,topsep=4pt,parsep=0pt}
\setlist[enumerate]{leftmargin=1.8em,itemsep=2pt,topsep=4pt,parsep=0pt}
\usepackage{authblk}
\usepackage{tcolorbox}
\tcbuselibrary{skins,breakable}
\tcbset{
  abstractstyle/.style={enhanced,breakable,colframe=natblue,colback=lightblue!50,
    boxrule=1.2pt,arc=5pt,top=10pt,bottom=10pt,left=12pt,right=12pt,
    fonttitle=\large\bfseries\color{natblue},
    attach boxed title to top center={yshift=-3mm},
    boxed title style={enhanced,colback=white,colframe=natblue,arc=3pt,boxrule=0.8pt}},
  hypothesisstyle/.style={enhanced,breakable,colframe=natgreen,colback=lightgreen!60,
    boxrule=1pt,arc=4pt,top=8pt,bottom=8pt,left=10pt,right=10pt},
  problemstyle/.style={enhanced,breakable,colframe=natorange,colback=lightyellow!80,
    boxrule=1pt,arc=4pt,top=8pt,bottom=8pt,left=10pt,right=10pt},
  outcomestyle/.style={enhanced,breakable,colframe=natblue!60,colback=lightblue!30,
    boxrule=0.8pt,arc=3pt,top=6pt,bottom=6pt,left=8pt,right=8pt}
}
\usepackage[colorlinks=true,linkcolor=natblue,citecolor=natblue,urlcolor=natblue2,
  pdftitle={端到端纠错码AI全自动设计框架}]{hyperref}
\usepackage{fancyhdr}
\pagestyle{fancy}
\fancyhf{}
\renewcommand{\headrulewidth}{0.5pt}
\renewcommand{\footrulewidth}{0pt}
\fancyhead[L]{\small\color{natgray}端到端纠错码AI全自动设计框架}
\fancyhead[R]{\small\color{natgray}\thepage}
\fancyfoot[C]{\small\color{natgray}\textit{Nature 投稿提案~$\cdot$~机密}}
\usepackage{caption}
\captionsetup{labelfont={bf,color=natblue},font=small,skip=5pt,labelsep=period}
\usepackage{setspace}
\setstretch{1.35}
\setlength{\parskip}{4pt}
\setlength{\parindent}{2em}
\theoremstyle{definition}
\newtheorem{hypothesis}{假说}[section]
\newtheorem{sciprob}{科学问题}[section]
\newtheorem{remark}{注记}[section]
\newcommand{\zstruct}{\bm{z}_{\mathrm{struct}}}
\newcommand{\zchannel}{\bm{z}_{\mathrm{channel}}}
\newcommand{\SNR}{\mathrm{SNR}}
\newcommand{\BER}{\mathrm{BER}}
\newcommand{\FER}{\mathrm{FER}}
\newcommand{\R}{\mathbb{R}}
\newcommand{\E}{\mathbb{E}}
\newcommand{\calT}{\mathcal{T}}
\newcommand{\calC}{\mathcal{C}}
\newcommand{\calF}{\mathcal{F}}
\newcommand{\bvec}[1]{\bm{#1}}
\newcommand{\thetacode}{\bm{\theta}_{\mathrm{code}}}
\newcommand{\thetahw}{\bm{\theta}_{\mathrm{hw}}}
\usepackage{microtype}
\title{
  {\color{natblue}\rule{\linewidth}{2.5pt}}\\[8pt]
  {\LARGE\bfseries\color{natblue}端到端纠错码AI全自动设计框架}\\[5pt]
  {\large\bfseries\color{natblue!80!black}从模糊需求到RTL代码的统一神经算子优化范式}\\[7pt]
  {\color{natblue!30}\rule{\linewidth}{0.6pt}}\\[4pt]
  {\normalsize\itshape\color{natgray}
    A Universal Neural Operator Framework for\\[2pt]
    End-to-End Error-Correcting Code Algorithm--Hardware Co-Design}\\[5pt]
  {\color{natblue!20}\rule{\linewidth}{0.4pt}}\\[-6pt]
}
\author[1]{\large [第一作者]}
\author[2]{\large [通讯作者\textsuperscript{*}]}
\affil[1]{\normalsize\color{natgray}[所在单位一]，[城市]}
\affil[2]{\normalsize\color{natgray}[所在单位二]，[城市]}
\date{\normalsize\color{natgray}\textit{研究提案}\quad$\cdot$\quad 2026年4月}
\begin{document}
\maketitle
\thispagestyle{empty}
\vspace{-6pt}
{\color{natgray!50}\hrule height 0.3pt}
\vspace{6pt}
\begin{tcolorbox}[abstractstyle,
  title={执行摘要 \quad {\normalsize\normalfont\itshape (Executive Summary)}}]

纠错码（Error-Correcting Code, ECC）是贯穿现代信息社会全域的核心基础技术——无论是5G/6G无线通信、相干光纤传输、NAND Flash固态存储、量子计算系统，还是新型存算一体器件，无一例外地高度依赖纠错码保障信息可靠传输与处理。然而，尽管应用场景千差万别，当前ECC设计范式却被三大结构性困境所困：\textbf{（一）过度依赖专家经验}——从码型构造、度分布优化到译码算法选择，全程仰赖跨越信息论、信号处理与VLSI设计的深厚专家知识壁垒；\textbf{（二）优化严重依赖大规模仿真}——获取一条可靠的BER性能曲线须耗费数小时乃至数天的蒙特卡洛仿真，使大规模码型搜索在工程上几乎不可行；\textbf{（三）算法优化与硬件实现孤立割裂}——译码算法与芯片硬件设计两条并行轨道互不知情，导致算法确定后硬件端大量推倒重来，严重拖慢标准化进程。

\vspace{4pt}
本提案提出 \textbf{PrimeCoDec}——一个面向信息系统纠错码需求的\textbf{全链路AI自动化框架}（Full-Pipeline AI Co-Design for Coding and Decoding）。PrimeCoDec 将上述三大困境同时消解：用户仅需以自然语言输入原始需求，框架即可\textbf{全自动生成}从纠错码编码方案、译码算法、硬件部署方案直至行业标准草案的完整输出，全程无需专家干预。

\vspace{4pt}
框架的核心创新体现于四个均衡配置的技术层次：\textbf{（一）编码方案AI自动化}，以AI驱动的码型搜索替代专家主导的码结构设计，实现面向任意物理信道的最优码型自动生成；\textbf{（二）无仿真译码算法优化}，以DeepONet神经算子在极低延迟内预测完整BER/FER性能曲线，从根本上替代耗时的蒙特卡洛仿真，实现译码算法参数的快速自动调优；\textbf{（三）算法--硬件协同自动实现}，以联合代理模型在统一目标函数下同步优化算法性能与VLSI指标，支持共模架构（一套硬件实现多种码型译码）等高价值设计输出；\textbf{（四）顶层需求形式化与标准自动输出}，以大语言模型（LLM）将自然语言需求转化为形式化约束，并将优化结果逆向映射为符合3GPP/IEEE格式的行业标准草案。

\vspace{4pt}
作为核心验证场景，我们选择\textbf{5G NR标准（3GPP Release 17）}作为基准——这是当前工业界极少数经数十年专家深度打磨的完备行业标准之一，对比天然公平且极具说服力。PrimeCoDec将全自动生成Polar码与LDPC码的最优方案，给出当前场景下的最优译码算法，设计共模架构译码器（一套硬件同时实现Polar码与LDPC码译码），并输出修订后的3GPP编码标准草案——从用户输入到可部署标准，全程由框架自动完成。
\end{tcolorbox}
\vspace{4pt}
\newpage
{\hypersetup{hidelinks}\setstretch{1.1}\tableofcontents}
\setstretch{1.35}
\newpage
\section{研究背景与核心动机}\label{sec:background}

\begin{tcolorbox}[hypothesisstyle,
  title={\normalsize\bfseries\color{natgreen}PrimeCoDec 故事脉络 \quad {\normalfont\itshape Story Narrative}}]
\textbf{一、问题根源}——纠错码在无线通信、光通信、数据存储、量子计算、存算一体等\textbf{所有信息系统}中不可或缺，但当前设计范式仍停留在``专家手工作坊''模式：过度依赖专家经验、优化严重依赖大规模仿真、算法与硬件深度割裂——三大困境导致ECC设计高度低效，严重制约新兴信息系统的快速部署。

\vspace{6pt}
\textbf{二、框架主张}——\textbf{PrimeCoDec} 提出：用户仅输入原始需求，框架即可全自动完成从纠错码编码设计、译码算法选择、硬件实现到行业标准输出的全链路自动化，打破ECC设计中一切专家壁垒。

\vspace{6pt}
\textbf{三、技术支柱}——四个均衡的技术层次共同支撑框架：
\begin{itemize}[nosep]
  \item \textbf{编码层}：AI驱动的码型自动搜索与构造（跨域通用信道表示）
  \item \textbf{译码层}：神经算子替代蒙特卡洛仿真，实现快速算法评估与调优
  \item \textbf{硬件层}：算法-硬件联合代理优化，支持共模架构自动实现
  \item \textbf{标准层}：LLM驱动的需求形式化与行业标准草案自动生成
\end{itemize}

\vspace{6pt}
\textbf{四、核心验证}——以\textbf{5G NR标准（3GPP Release 17）}为基准，验证PrimeCoDec的全链路能力：从自然语言需求输入，到Polar码/LDPC码最优方案、共模译码器硬件，最终输出修订后的3GPP标准草案——一个完整的从需求到标准的自动化闭环。
\end{tcolorbox}

\vspace{8pt}
\subsection{纠错码：横跨全域信息系统的核心基础技术}

纠错码技术诞生于1948年Shannon信息论，历经七十余年发展，已成为数字信息基础设施不可或缺的核心组件——其触角延伸至当今几乎所有信息系统，是横跨无线通信、光通信、数据存储、新型计算乃至量子信息的\textbf{共性底层技术}。

\textbf{无线通信}方面，在5G NR标准（3GPP Release 15+）中，LDPC码\cite{gallager1962}承担eMBB场景数据信道编码任务，Polar码\cite{arikan2009}被采纳为控制信道方案。下一代6G系统对URLLC（超可靠低时延通信）提出了更严苛的编码需求——在更短码长（$n < 256$）与更高可靠性（BLER $\leq 10^{-9}$）约束下保障极端低时延。

\textbf{光通信}方面，软判决前向纠错（SD-FEC）已广泛采用LDPC码，实现净编码增益（NCG）超过10~dB，成为400G/800G乃至未来Tb/s光传输系统的关键使能技术。

\textbf{存储系统}方面，随着NAND Flash存储密度从SLC向QLC跃升，信道特性呈现显著的非高斯性与时变性，亟需针对磨损非平稳噪声的定制化ECC方案。

\textbf{量子计算}方面，量子纠错码（QECC）是实现容错量子计算的唯一已知途径\cite{shor1995}。以表面码\cite{fowler2012}为代表的拓扑量子纠错方案高度依赖对量子噪声模型的精确建模与高效经典译码。

\textbf{存算一体}方面，忆阻器（Memristor）阵列的模拟电导分布噪声构成独特的存储信道模型，现有ECC体系缺乏针对此类新型器件的系统性设计工具。

上述五大应用领域共同指向一个深刻事实：\textbf{只要存在信息的传输与存储，纠错码就无处不在}。然而令人震惊的是，面向如此广泛且多样化的应用，当前ECC的设计方式却几乎停留在香农时代的手工作坊模式。

\subsection{现有ECC设计范式的三大结构性困境}\label{sec:pain_points}

纵观上述所有应用场景，ECC设计者每次面对新需求时，都不得不从头开始应对同样的三大系统性困境：

\textbf{困境一：码型优化与算法设计高度依赖专家先验知识。}LDPC基图设计需要深厚的环分析、度分布优化与密度进化理论基础\cite{richardson2001}；Polar码编码序列的选择依赖Tal--Vardy算法\cite{tal2013}等复杂理论工具；译码算法的参数调优（迭代次数、量化位宽、调度策略）需要多年工程经验积累。这一\textbf{专家壁垒}严重制约了ECC技术向新型应用场景的快速推广，也使得每当出现新系统、新信道时，ECC方案的从头设计成为不可避免的高成本工程。

\textbf{困境二：性能评估与参数优化严重依赖大规模仿真。}获取一条可靠的BER曲线需要在每个目标SNR点累计数亿次随机噪声样本。以码长$n = 10^4$、目标BER为$10^{-8}$的LDPC码为例，完整评估需约$10^{12}$次信道仿真，即便借助高端GPU并行加速也需数小时乃至数天。这一瓶颈使得大规模码型参数搜索几乎不可行，\textbf{大量人力与计算资源被耗费在重复仿真而非真正的创造性设计上}。

\textbf{困境三：算法优化与硬件实现深度脱节。}算法工程师在优化BER性能时通常不考虑译码器的硬件实现代价；当算法确定后转入硬件实现阶段，往往发现在面积/功耗约束下无法直接实现，导致大量推倒重来。更深层的问题是：\textbf{当前没有任何框架能够将算法优化与硬件部署、乃至行业标准制定统一于同一闭环}，每个环节都是孤立的"信息孤岛"。

这三大困境共同造成了ECC领域高度低效的设计生态：仅有极少数行业（如5G NR无线通信）凭借数十年专家深度投入，才建立起完备且充分优化的行业标准（3GPP TS 38.212等\cite{3gpp38212}）；绝大多数新兴应用场景（量子计算、忆阻器存算一体、6G新波形等）至今仍缺乏系统性的ECC设计支撑。

\subsection{人工智能带来的变革机遇}\label{sec:ai_opportunity}

近年来，AI技术在若干关键方向取得的突破为克服上述三大困境、实现ECC设计全链路自动化提供了历史性机遇：

\textbf{图神经网络（GNN）用于结构化学习。}因子图是LDPC码、Polar码等现代纠错码的统一数学表示工具，GNN天然适合对因子图的拓扑结构进行编码学习，已有大量工作证明其在神经译码\cite{nachmani2018,kang2022}领域的有效性。

\textbf{神经算子（Neural Operator）用于泛函逼近。}以DeepONet\cite{lu2021}与FNO\cite{kovachki2023}为代表的神经算子，已在PDE求解、湍流预测等领域展现出强大的泛化能力，为脱离蒙特卡洛仿真提供了全新理论工具。

\textbf{大语言模型（LLM）用于需求理解与代码生成。}GPT-4o、Claude等LLM已展示出强大的自然语言理解与代码生成能力，在EDA领域的初步应用\cite{liu2023,chang2024}表明RTL代码生成与硬件需求理解的可行性。

\textbf{硬件代理模型（Hardware Surrogate Models）。}以GNN与MLP为核心的轻量硬件性能代理模型\cite{ustun2020,geng2021}在逻辑综合结果预测方面展现出接近综合工具精度的同时实现了数量级的速度提升。这为算法与硬件的协同优化提供了现实可行的计算基础。

上述四个技术方向与ECC设计的三大困境构成精准映射：GNN解决码结构学习的专家依赖，神经算子突破仿真瓶颈，LLM打通需求与标准之间的语义鸿沟，硬件代理模型弥合算法与硬件之间的割裂。\textbf{PrimeCoDec}正是将这四大技术有机融合，构建覆盖ECC设计全链路的统一自动化框架的系统性尝试。
\section{核心科学问题}\label{sec:problems}

本研究的学术价值根植于若干具有根本意义的开放科学问题。围绕PrimeCoDec框架所要解决的三大结构性困境，我们将其提炼为以下五个相互关联的核心科学问题：

\vspace{6pt}
\begin{tcolorbox}[problemstyle]
\textbf{科学问题 P1：通用信道统一表示}

如何为来自截然不同物理机制（电磁传播、光子噪声、量子退相干、固态器件噪声）的异质噪声信道，建立一种\textbf{数学完备且计算高效}的统一表示框架，使得基于此框架设计的AI模型能够在不同系统间无缝迁移，而无需针对每类系统重新训练或重新设计？

\textit{挑战所在}：不同物理系统的噪声统计结构差异极大——从高斯分布（AWGN）、混合高斯（Flash）到Pauli噪声（量子）；噪声的时空相关性结构从无关（独立信道）到强相关（突发错误信道）跨越多个数量级；概率密度函数可能不存在解析形式（仅有测量数据）。
\end{tcolorbox}

\vspace{6pt}
\begin{tcolorbox}[problemstyle]
\textbf{科学问题 P2：纠错码性能的高精度神经算子预测}

是否存在一种神经网络架构，能够在\textbf{无需蒙特卡洛仿真}的前提下，对任意给定的码结构$\calC$与信道三元组$\calT$，以\textbf{高精度}（BER曲线误差$< 0.1$~dB）、\textbf{超快速度}（推理时间$< 10$~ms）预测完整的BER/FER性能曲线，并对不同码长、码率、校验矩阵与信道参数的组合具有强泛化能力？

\textit{挑战所在}：该映射本质上是从无限维函数空间到无限维函数空间的算子，传统神经网络仅能映射有限维向量，直接应用存在根本性局限。BER曲线在瀑布区（waterfall region）呈现高度非线性的急降特性，对预测模型的局部精度提出了极高要求；BER值跨越$10^{-1}$到$10^{-10}$多个数量级，需要特殊的损失函数设计。
\end{tcolorbox}

\vspace{6pt}
\begin{tcolorbox}[problemstyle]
\textbf{科学问题 P3：算法-硬件联合性能代理建模}

如何构建一个轻量级代理模型，能够同时以高精度预测ECC\textbf{算法层指标}（BER/FER性能）与\textbf{硬件层指标}（芯片面积$A$、功耗$P$、关键路径时延$T_d$、吞吐量$\Gamma$），从而支持在统一目标函数框架下的算法-硬件协同优化？

\textit{挑战所在}：算法性能与硬件指标的影响因素部分重叠、部分正交。硬件指标高度依赖工艺节点与综合工具设置，如何在代理模型中有效编码工艺参数是一个未解问题。
\end{tcolorbox}

\vspace{6pt}
\begin{tcolorbox}[problemstyle]
\textbf{科学问题 P4：大语言模型驱动的需求形式化}

如何设计一个可靠的LLM推理管道，将工程师提出的\textbf{自然语言模糊需求}（如"需要一个用于5G URLLC的、在城市非视距（NLOS）场景下可靠性达$10^{-9}$的纠错方案"）或标准规范文档中的\textbf{上下文约束}（3GPP TS 38.211等），自动转化为优化引擎可直接处理的\textbf{形式化约束集合}$\calC$？该过程的\textbf{完备性}与\textbf{一致性}如何得到保证？

\textit{挑战所在}：工程需求往往包含隐含假设、领域专有术语与相互竞争的多目标约束。LLM在处理此类结构化工程推理时容易出现幻觉（hallucination）或约束遗漏，需要配套的验证机制。
\end{tcolorbox}

\vspace{6pt}
\begin{tcolorbox}[problemstyle]
\textbf{科学问题 P5：行业标准自动化生成}

如何将PrimeCoDec框架输出的最优码型参数、译码算法与硬件实现结果，\textbf{自动逆向映射}为符合3GPP、IEEE 802.11等工业格式的行业标准草案文档？该自动化标准生成过程的\textbf{合规性}、\textbf{完备性}与\textbf{可读性}如何保证？

\textit{挑战所在}：行业标准文档具有严格的格式规范、术语体系与条款逻辑；将数学优化结果翻译为工程规范语言需要深度领域知识；不同应用场景（无线、光通信、存储）对应不同的标准化组织与文档体系，如何实现统一的标准生成接口是一个未解问题。
\end{tcolorbox}

\vspace{6pt}
上述五个科学问题相互关联、层层递进，共同构成PrimeCoDec的理论核心。P1是整个框架的数学基础，使跨域泛化成为可能；P2是突破仿真瓶颈、消解专家依赖的核心技术；P3将性能预测能力扩展至硬件维度，终结算法-硬件割裂；P4是实现端到端自动化闭环的用户交互接口；P5则将框架的价值延伸至行业标准层，使优化结果直接转化为可落地的工业规范。本提案将系统性地攻克这五个问题，建立完整的理论与工程框架。
\section{科学假说与研究目标}\label{sec:hypotheses}

基于对上述四个核心科学问题的分析，我们提出以下四个可证伪的科学假说，并据此确立具体的量化研究目标：

\subsection{科学假说}

\begin{tcolorbox}[hypothesisstyle]
\textbf{假说 H1（通用信道抽象的完备性）}

任意物理系统中的信道噪声均可由信道三元组
$\calT = \bigl(\{\SNR_i\}_{i=1}^n,\;\bvec{R} \in \R^{n\times n},\;p(x)\bigr)$
精确刻画。基于此三元组表示训练的ECC性能预测模型，在未见过的物理系统（仅提供对应三元组）上的预测误差，与在同类系统上训练的专用模型相当。

\textbf{验证方案}：在五类物理系统上独立获取信道三元组，使用统一模型预测，对比专用模型精度，目标：通用模型精度损失$< 0.2$~dB。
\end{tcolorbox}

\vspace{4pt}
\begin{tcolorbox}[hypothesisstyle]
\textbf{假说 H2（神经算子逼近的有效性）}

从（码结构，信道三元组）到BER/FER性能曲线的映射$\calF: (\calC, \calT) \mapsto \BER(\cdot)$构成一个定义在适当Banach空间上的连续算子，可由DeepONet以任意精度逼近（由DeepONet的万能逼近定理保证~\cite{lu2021}）。具体而言，存在有限参数量的DeepONet模型，能够在给定精度$\epsilon = 0.1$~dB下预测任意LDPC码与Polar码在任意物理信道三元组上的完整BER/FER曲线。

\textbf{验证方案}：构建包含$> 10^6$个（码, 信道, BER曲线）三元组的训练集，训练DeepONet，在独立测试集上验证BER预测RMSE（dB域）$< 0.1$~dB，曲线整体相对误差$< 5\%$。
\end{tcolorbox}

\vspace{4pt}
\begin{tcolorbox}[hypothesisstyle]
\textbf{假说 H3（硬件性能的低维可学习性）}

LDPC/Polar译码器的硬件实现指标（逻辑面积$A$、动态功耗$P$、关键路径时延$T_d$）是码参数与工艺参数的低维可学习函数，可由轻量GNN+MLP代理模型以$< 5\%$相对误差精确预测，同时推理速度较RTL综合工具快$> 100$倍。

\textbf{验证方案}：使用工业综合工具（Design Compiler）生成$> 5000$个设计数据点，训练代理模型，验证面积/功耗预测误差$< 5\%$，时延误差$< 8\%$。
\end{tcolorbox}

\vspace{4pt}
\begin{tcolorbox}[hypothesisstyle]
\textbf{假说 H4（端到端全自动化的可行性）}

通过LLM需求解析引擎、神经算子性能预测器与硬件代理模型、以及参数化RTL生成器的有机组合，可以实现从自然语言需求输入到可综合RTL代码输出的\textbf{全自动闭环设计流水线}，且最终RTL代码在FPGA原型验证中的实测BER性能与代理模型预测值之间的误差$< 0.3$~dB。

\textbf{验证方案}：以5种不同应用场景的自然语言需求为输入，运行完整E3A流水线，最终RTL代码在Xilinx Alveo FPGA平台进行硬件在环（HIL）验证。
\end{tcolorbox}

\subsection{量化研究目标}\label{sec:objectives}

基于上述假说，本研究设定以下可量化的研究目标，作为项目验收标准：

\begin{table}[htbp]
\centering
\caption{E3A框架量化研究目标}
\label{tab:objectives}
\rowcolors{2}{lightblue!30}{white}
\begin{tabular}{@{}L{3.5cm}L{5.0cm}L{3.5cm}@{}}
\toprule
\textbf{目标维度} & \textbf{具体指标} & \textbf{目标值} \\
\midrule
性能预测精度 & BER曲线预测误差（dB域RMSE） & $< 0.1$~dB \\
性能预测速度 & 单次BER曲线预测时间 & $< 10$~ms \\
相对仿真加速比 & vs.\ GPU加速蒙特卡洛 & $> 1000\times$ \\
硬件面积预测精度 & 逻辑面积相对误差 & $< 5\%$ \\
硬件功耗预测精度 & 动态功耗相对误差 & $< 5\%$ \\
端到端优化速度 & 从需求到最优参数（含100次迭代） & $< 10$~min \\
RTL验证一致性 & FPGA实测vs.预测BER误差 & $< 0.3$~dB \\
系统通用性 & 成功验证的物理系统类型数 & $\geq 5$ \\
需求形式化准确率 & 关键约束提取准确率 & $> 95\%$ \\
\bottomrule
\end{tabular}
\end{table}

\subsection{研究范围说明}

本研究聚焦于以下类型的纠错码：
\begin{itemize}
  \item \textbf{LDPC码}：包括3GPP NR标准基图（BG1/BG2）及自定义扩展基图
  \item \textbf{Polar码}：包括标准CRC辅助Polar码（CA-Polar）及PAC码
\end{itemize}

目标物理系统涵盖：5G无线信道、相干光通信信道、NAND Flash存储信道、量子比特噪声信道（超导体系）及忆阻器存储信道（第\ref{sec:applications}节详述）。
\section{科学贡献与应用前景}\label{sec:contribution}

\subsection{四大核心创新点}

PrimeCoDec的创新贡献覆盖ECC全设计链路的四个关键层次，各层次均具备独立的科学意义与工程价值，同时在统一框架内深度耦合协同：

\vspace{4pt}
\textbf{创新点一：编码方案AI自动化（Encoding Layer）}

针对当前码型设计高度依赖专家先验知识的困境，PrimeCoDec以\textbf{通用信道三元组}$\calT{=}(\{\SNR_i\},\bvec{R},p(x))$统一表征任意物理信道，并以AI驱动的码型搜索替代专家主导的手工设计。基于因子图的图神经网络编码码结构拓扑，以进化算法与贝叶斯优化在统一信道表示下自动探索LDPC基图与Polar码冻结比特序列的最优配置。这一创新首次使ECC编码方案设计对任意物理信道具备\textbf{领域无关的自动化能力}，无需针对每类系统进行专家重新设计。

\vspace{4pt}
\textbf{创新点二：无仿真译码算法自动优化（Decoding Layer）}

针对性能评估严重依赖耗时蒙特卡洛仿真的困境，PrimeCoDec以\textbf{神经算子（DeepONet）}建立从码结构-信道三元组到完整BER/FER性能曲线的泛函逼近器。Branch网络以图变换器编码码字的因子图拓扑，Trunk网络以三路并行编码器处理信道三元组，二者通过神经内积融合，以极低延迟替代蒙特卡洛仿真，实现对任意码长、码率、译码算法配置（迭代次数、量化位宽、调度策略等）的快速性能预测与自动调优。这一创新从根本上\textbf{破除了ECC优化对大规模仿真的依赖}，将评估-优化闭环从天级压缩至秒级。

\vspace{4pt}
\textbf{创新点三：算法--硬件协同自动实现（Hardware Layer）}

针对算法优化与硬件实现深度脱节的困境，PrimeCoDec构建\textbf{硬件性能代理模型}（GNN+MLP），建立从码参数、译码算法参数到VLSI实现指标（面积、功耗、时延、吞吐量）的代理映射，并与算法性能代理在统一目标函数下进行\textbf{联合帕累托优化}。框架原生支持\textbf{共模架构}设计——以一套硬件同时实现多种码型（如Polar码与LDPC码）译码，通过联合代理模型在设计空间中自动定位共模最优点。最终由参数化RTL模板库结合LLM代码补全，输出可综合的System Verilog译码器代码。

\vspace{4pt}
\textbf{创新点四：顶层需求形式化与行业标准自动输出（Standard Layer）}

针对ECC设计闭环缺失顶层输入与标准化输出接口的问题，PrimeCoDec通过两端创新实现从用户到标准的完整通路：\textbf{输入端}，大语言模型（LLM）以RAG知识检索与CoT推理将自然语言模糊需求（或3GPP/IEEE规范文档）自动转化为形式化约束集合，并进行香农极限可行性验证；\textbf{输出端}，框架将最优码型参数、译码算法配置与硬件实现结果自动逆向映射为符合3GPP/IEEE/ITU-T格式的\textbf{行业标准草案文档}，支持多场景标准配置集的批量生成。这一创新首次使ECC设计闭环延伸至标准制定层面，实现真正意义上的``需求到标准''全自动化。

\subsection{PrimeCoDec能力矩阵}

\begin{table}[htbp]
\centering
\caption{PrimeCoDec框架四层能力矩阵}
\label{tab:capability}
\setlength{\tabcolsep}{4pt}
\rowcolors{2}{lightblue!15}{white}
\begin{tabular}{@{}L{2.8cm}L{3.5cm}L{3.5cm}L{3.0cm}@{}}
\toprule
\textbf{技术层次} & \textbf{当前范式局限} & \textbf{PrimeCoDec方式} & \textbf{科学意义} \\
\midrule
编码设计 & 专家手工设计，高度依赖领域知识 & AI驱动码型搜索，通用信道三元组抽象 & 编码设计民主化 \\
译码算法 & Monte Carlo仿真，数小时乃至数天 & 神经算子性能预测，极低延迟 & 仿真瓶颈的根本破除 \\
硬件实现 & 算法-硬件孤立迭代，大量返工 & 联合代理优化，共模架构自动生成 & 首个算法-硬件协同ECC工具 \\
标准输出 & 无自动化工具，需专家人工撰写 & LLM形式化+自动标准草案生成 & ECC设计延伸至标准化层 \\
\bottomrule
\end{tabular}
\end{table}

\subsection{核心验证场景：以5G NR为基准的完整闭环}\label{sec:5gnr_case}

当前工业界，\textbf{仅有极少数领域}具备经过数十年专家深度打磨的完备行业标准——5G NR无线通信标准（3GPP Release 17\cite{3gpp38212}）正是其中的典范，代表了全球顶尖ECC专家数十年集体智慧的结晶。以此为基准，可以最直观地展示PrimeCoDec框架从需求到标准的全链路能力，且该对比天然公平：5G NR标准已被充分优化，没有``未经优化''的先天劣势。

\textbf{场景输入}：以自然语言描述5G NR eMBB数据信道的典型场景规格（AWGN信道，目标BER约束、延迟约束、芯片面积与吞吐量约束，与3GPP TS 38.212 Release 17规定的参考场景对齐）。PrimeCoDec框架全自动运行，无人工干预。

\textbf{四层全链路输出}：
\begin{itemize}
  \item \textbf{编码层输出}：针对5G NR eMBB场景，PrimeCoDec自动生成Polar码（控制信道）与LDPC码（数据信道）的最优码型参数配置，包括基图结构、码长码率组合与冻结比特序列。
  \item \textbf{译码层输出}：框架以神经算子扫描全部算法配置空间，自动确定当前场景下的最优译码算法——Polar码采用SCL（Successive Cancellation List）译码，LDPC码采用分层置信传播（Layered BP）译码，并给出最优迭代次数、量化位宽等参数。
  \item \textbf{硬件层输出}：PrimeCoDec自动生成\textbf{共模架构译码器}——一套硬件平台同时支持Polar码与LDPC码译码（根据控制信道/数据信道需求自动切换），并输出可综合的System Verilog RTL代码与综合脚本。
  \item \textbf{标准层输出}：框架自动生成符合3GPP文档格式的修订标准草案，包含新的编码序列表、基图参数集与译码器配置建议，形成可直接提交标准委员会审阅的文档。
\end{itemize}

这一案例清晰勾勒出PrimeCoDec作为``ECC全链路自动化引擎''的完整能力边界：接受自然语言需求，输出可部署的行业标准，中间全部环节由AI自动完成。

\subsection{跨领域应用前景}\label{sec:applications}

\begin{table}[htbp]
\centering
\caption{PrimeCoDec框架在五大物理系统中的应用价值}
\label{tab:applications}
\setlength{\tabcolsep}{4pt}
\rowcolors{2}{lightblue!15}{white}
\begin{tabular}{@{}L{2.0cm}L{3.2cm}L{3.2cm}L{4.8cm}@{}}
\toprule
\textbf{应用系统} & \textbf{当前痛点} & \textbf{PrimeCoDec解决方式} & \textbf{核心价值} \\
\midrule
5G/6G无线 & 标准更迭频繁，设计周期长 & 需求直接输入，自动生成最优方案及3GPP标准草案 & 根本压缩设计周期，使能快速标准化 \\
相干光通信 & 非高斯噪声难建模 & 非参数PDF编码器天然适配 & 突破非高斯信道ECC性能瓶颈 \\
NAND Flash & 磨损噪声分布非平稳 & 逐单元SNR序列精确建模 & 定制化ECC延长存储寿命 \\
量子纠错 & 量子码设计无EDA工具 & 三元组适配Pauli信道 & 开辟量子EDA新方向 \\
忆阻器存算 & 无信道先验模型 & 数据驱动PDF学习 & 使存算一体器件具备系统性ECC支撑 \\
\bottomrule
\end{tabular}
\end{table}

\textbf{典型新兴案例：忆阻器存算一体ECC设计}。忆阻器阵列目前\textbf{没有任何成熟的ECC适配工具}，现有方案依赖纯经验调试。PrimeCoDec仅凭实测噪声样本（核密度估计构建$p(x)$）即可完成设计：将器件特性封装为三元组（非均匀SNR序列+行列线串扰相关性矩阵+非参数PDF），神经算子预测该新型信道下LDPC码的性能曲线，自动输出适配该信道的最优码型、译码算法与RTL代码——无需预先存在的信道模型，无需领域专家介入。

\subsection{学术贡献与科学影响}\label{sec:impact}

\subsubsection{新范式的确立}

PrimeCoDec框架的核心科学贡献在于\textbf{确立了信息论与AI深度融合的全链路自动化新范式}：以神经算子理论为桥梁，将函数逼近论引入纠错码设计；以LLM为界面，将用户需求与行业标准直接打通；以算法-硬件联合代理为纽带，终结ECC设计中延续数十年的``串行孤岛''模式。

这一范式转变具有标志性意义：（1）将``码-信道''关系从``点仿真''提升为``算子学习''，赋予ECC设计连续函数空间的理论框架；（2）将硬件实现纳入ECC设计的一阶约束，终结``先算法后硬件''的串行模式；（3）将行业标准生成纳入设计闭环，使``从需求到标准''的全自动化成为现实；（4）通过LLM消除ECC专家知识壁垒，重新定义谁可以设计纠错码。

\subsubsection{信息论层面的数学贡献}

\textbf{统一信道表示}：三元组$(\{\SNR_i\},\bvec{R},p(x))$提供了一种新的信道数学语言，可能演化为信道类型不可知ECC分析的通用工具，推动信道容量、编码定理与极化现象的统一表述。

\textbf{算子视角的编码定理}：将ECC性能（BER/FER）理解为从码参数到SNR的连续算子，为利用函数分析工具研究编码增益、极化现象与容量可达性开辟新途径。

\subsubsection{跨学科科学影响}

\begin{itemize}[nosep,leftmargin=1.5em]
  \item \textbf{AI for Science}：PrimeCoDec是``AI加速科学发现''（AI4Science）范式在信息论领域的首个系统性实践，与AlphaFold（蛋白质结构）、AI for Fluid Dynamics（CFD算子学习）的范式相似，有望引发信息论界的范式效应。
  \item \textbf{EDA方法论}：将机器学习代理模型引入ECC硬件设计空间探索，与EDA领域的``机器学习增强设计自动化''趋势深度对接，为VLSI设计方法论提供新的方向。
  \item \textbf{量子计算}：为量子纠错码设计提供首个AI辅助工具，可能成为容错量子计算机工程化路径中的重要基础设施。
  \item \textbf{新型存储器}：为忆阻器、MRAM等新型存储技术提供无先验信道模型的ECC设计能力，赋能下一代存储芯片的可靠性工程。
  \item \textbf{标准化加速}：PrimeCoDec的标准生成模块有潜力成为下一代无线、光通信、量子通信标准（6G、400G+、量子网络）制定过程中的AI辅助工具，大幅压缩标准化周期。
\end{itemize}

\subsubsection{开放科学承诺}

本项目承诺全面开放：\textbf{代码}（GitHub Apache-2.0开源）、\textbf{训练数据}（Zenodo数据仓库，含大规模DC综合结果与性能预测标签）、\textbf{预训练模型权重}（Hugging Face模型库）、\textbf{基准测试框架}（标准化的五信道类型评估协议）、\textbf{标准生成模板库}（覆盖3GPP/IEEE/ITU-T格式），为后续研究者提供完整可复现的研究基础。

\section{技术可行性论证}\label{sec:feasibility}

\subsection{神经算子理论保证}

E3A框架的核心算法性能预测模块基于DeepONet，其数学基础为Chen \& Chen（1995）关于神经算子的万能逼近定理\cite{chen1995}：

\begin{tcolorbox}[hypothesisstyle,title={定理：神经算子万能逼近性（Chen \& Chen, 1995）}]
设$\calF: C(K_1;\R^d) \to C(K_2;\R)$为定义在紧集上的连续非线性算子，$\sigma$为非多项式连续激活函数。则对任意$\varepsilon > 0$，存在正整数$p$与神经网络参数，使得
$$
\sup_{u \in C(K_1;\R^d)} \left|\calF(u)(y) - \sum_{i=1}^{p} b_i(u) \cdot t_i(y)\right| < \varepsilon, \quad \forall y \in K_2.
$$
\end{tcolorbox}

在ECC性能预测中，$u$对应信道三元组$\calT$（连续函数），$y$对应SNR评估点，$\calF$对应BER/FER性能算子。定理保证框架在理论上具备无限逼近精度，与码型和信道类型无关。

\subsection{DeepONet在通信系统中的先验验证}

Lu等（2021）\cite{lu2021}在\textit{Nature Machine Intelligence}上首次报告DeepONet在物理场算子学习中的实验验证，证明神经算子在复杂非线性算子逼近中兼具高精度（误差$<1\%$）与高效推理（加速比$>10^4\times$）。本框架将此范式首次系统化迁移至ECC性能预测，理论基础充分，先前方法验证充分。

\subsection{已有相关工作支撑}

\begin{table}[htbp]
\centering
\caption{E3A各技术模块的先前工作支撑}
\label{tab:prior_work}
\setlength{\tabcolsep}{4pt}
\rowcolors{2}{lightblue!15}{white}
\begin{tabular}{@{}L{2.5cm}L{4.5cm}L{4.5cm}L{2.0cm}@{}}
\toprule
\textbf{技术模块} & \textbf{代表性先前工作} & \textbf{关键进展} & \textbf{本工作突破} \\
\midrule
因子图GNN & Nachmani等（2018）\cite{nachmani2018} & GNN显著优于BP & 推理不译码 \\
图变换器 & Dwivedi \& Bresson（2021）\cite{dwivedi2021} & 图注意力+位置编码 & 迁移至ECC图 \\
神经算子 & Lu等（2021）\cite{lu2021} & DeepONet理论与实验 & 首次用于ECC \\
硬件代理模型 & Ustun等（2020）\cite{ustun2020} & HLS性能机器学习预测 & 扩展至ECC \\
LLM for EDA & Chang等（2024）\cite{chang2024} & LLM辅助RTL生成 & 形式约束解析 \\
\bottomrule
\end{tabular}
\end{table}

\subsection{计算量估算}

\textbf{训练阶段}：性能预测代理模型单次训练约需$10^5$个（码参数，信道参数，性能标签）样本，可在$4\times$A100 GPU上于24小时内完成；硬件代理模型训练样本通过Design Compiler批量综合产生，单批次（100配置）约需8小时，并行化后可于1周内建库。\textbf{推理阶段}：单次神经算子推理耗时约$0.5$~ms（A100）或约$5$~ms（CPU），满足交互式设计需求。\textbf{与蒙特卡洛基准对比}：蒙特卡洛仿真达到BER$=10^{-9}$置信度需$\sim10^{11}$次仿真，在现代服务器上约需3--10小时；代理模型将其压缩至毫秒级，\textbf{加速比约$10^6\sim10^7\times$}。
\begin{thebibliography}{99}

\bibitem{gallager1962}
R.~G. Gallager,
``Low-density parity-check codes,''
\textit{IRE Trans. Inf. Theory}, vol.~8, no.~1, pp.~21--28, Jan.~1962.

\bibitem{arikan2009}
E.~Ar\i kan,
``Channel polarization: A method for constructing capacity-achieving codes for
symmetric binary-input memoryless channels,''
\textit{IEEE Trans. Inf. Theory}, vol.~55, no.~7, pp.~3051--3073, Jul.~2009.

\bibitem{shor1995}
P.~W. Shor,
``Scheme for reducing decoherence in quantum computer memory,''
\textit{Phys. Rev. A}, vol.~52, no.~4, pp.~R2493--R2496, 1995.

\bibitem{fowler2012}
A.~G. Fowler, M.~Martinis, J.~M. Martinis, and J.~M. Gambetta,
``Surface codes: Towards practical large-scale quantum computation,''
\textit{Phys. Rev. A}, vol.~86, no.~3, p.~032324, 2012.

\bibitem{richardson2001}
T.~J. Richardson and R.~L. Urbanke,
``The capacity of low-density parity-check codes under message-passing decoding,''
\textit{IEEE Trans. Inf. Theory}, vol.~47, no.~2, pp.~599--618, Feb.~2001.

\bibitem{tal2013}
I.~Tal and A.~Vardy,
``List decoding of polar codes,''
\textit{IEEE Trans. Inf. Theory}, vol.~61, no.~5, pp.~2213--2226, May~2015.

\bibitem{lu2021}
L.~Lu, P.~Jin, G.~Pang, Z.~Zhang, and G.~E. Karniadakis,
``Learning nonlinear operators via DeepONet based on the universal approximation
theorem of operators,''
\textit{Nat. Mach. Intell.}, vol.~3, no.~3, pp.~218--229, Mar.~2021.

\bibitem{kovachki2023}
N.~Kovachki \textit{et al.},
``Neural operator: Learning maps between function spaces with applications to PDEs,''
\textit{J. Mach. Learn. Res.}, vol.~24, no.~89, pp.~1--97, 2023.

\bibitem{chen1995}
T.~Chen and H.~Chen,
``Universal approximation to nonlinear operators by neural networks with arbitrary
activation functions and its application to dynamical systems,''
\textit{IEEE Trans. Neural Netw.}, vol.~6, no.~4, pp.~911--917, Jul.~1995.

\bibitem{nachmani2018}
E.~Nachmani, E.~Marciano, L.~Lugosch, W.~J. Gross, D.~Burshtein, and Y.~Be'ery,
``Deep learning methods for improved decoding of linear codes,''
\textit{IEEE J. Sel. Topics Signal Process.}, vol.~12, no.~1, pp.~119--131, 2018.

\bibitem{nachmani2019}
E.~Nachmani and Y.~Be'ery,
``Hyper-graph-network decoders for block codes,''
in \textit{Proc. NeurIPS}, 2019, pp.~2329--2339.

\bibitem{kang2022}
M.~Kang \textit{et al.},
``Polar code design for generalized polar decoding,''
\textit{IEEE Trans. Commun.}, vol.~70, no.~11, pp.~7108--7122, Nov.~2022.

\bibitem{liu2023}
S.~Liu, W.~Fang, Y.~Li, Y.~Zhuo, J.~Li, and T.~Shi,
``RTLCoder: Outperforming GPT-3.5 in design RTL generation with our open-source
dataset and lightweight solution,''
\textit{arXiv preprint}, arXiv:2312.08617, 2023.

\bibitem{chang2024}
C.-H. Chang, Y.-M. Yang, C.-C. Ho, Y.-K. Wang, and H.-Y. Chang,
``Data is all you need: Finetuning LLMs for chip design via an automated design-data
augmentation framework,''
in \textit{Proc. DAC}, 2024.

\bibitem{ustun2020}
E.~Ustun, C.~Deng, D.~Pal, Z.~Li, and Z.~Zhang,
``Accurate operation delay prediction for FPGA HLS using graph neural networks,''
in \textit{Proc. ICCAD}, 2020, pp.~1--9.

\bibitem{geng2021}
H.~Geng \textit{et al.},
``N-Parity: A machine learning approach for estimating power and area of a neural
network VLSI implementation,''
in \textit{Proc. DATE}, 2021.

\bibitem{ryan2009}
W.~E. Ryan and S.~Lin,
\textit{Channel Codes: Classical and Modern}.
Cambridge, U.K.: Cambridge Univ. Press, 2009.

\bibitem{lin2004}
S.~Lin and D.~J. Costello,
\textit{Error Control Coding}, 2nd~ed.
Upper Saddle River, NJ: Prentice-Hall, 2004.

\bibitem{3gpp38212}
3GPP,
``NR; multiplexing and channel coding (Release 17),''
3GPP TS 38.212 v17.3.0, Sep.~2022.

\bibitem{velickovic2018}
P.~Veli\v{c}kovi\'{c}, G.~Cucurull, A.~Casanova, A.~Romero, P.~Li\`{o}, and Y.~Bengio,
``Graph attention networks,''
in \textit{Proc. ICLR}, 2018.

\bibitem{dwivedi2021}
V.~P. Dwivedi and X.~Bresson,
``A generalization of transformers to graphs,''
\textit{arXiv preprint}, arXiv:2012.09699, 2021.

\bibitem{li2021fno}
Z.~Li, N.~Kovachki, K.~Azizzadenesheli, B.~Liu, K.~Bhattacharya, A.~Stuart, and A.~Anandkumar,
``Fourier neural operator for parametric partial differential equations,''
in \textit{Proc. ICLR}, 2021.

\bibitem{gruber2017}
T.~Gruber, S.~Cammerer, J.~Hoydis, and S.~ten Brink,
``On deep learning-based channel decoding,''
in \textit{Proc. CISS}, 2017.

\bibitem{cammerer2022}
S.~Cammerer, M.~Schmalen, L.~Schmalen, and S.~ten Brink,
``Learned factor graphs for inference from stationary time sequences,''
\textit{IEEE Trans. Signal Process.}, vol.~70, pp.~366--381, 2022.

\bibitem{zhao2023}
Z.~Zhao, D.~Fan, and Z.~Xu,
``Neural network-based turbo equalization for memristive channels,''
\textit{IEEE Trans. Commun.}, vol.~71, no.~4, pp.~2098--2112, Apr.~2023.

\bibitem{preskill2018}
J.~Preskill,
``Quantum computing in the NISQ era and beyond,''
\textit{Quantum}, vol.~2, p.~79, 2018.

\bibitem{gottesman1997}
D.~Gottesman,
``Stabilizer codes and quantum error correction,''
Ph.D. dissertation, Caltech, 1997.

\bibitem{wang2023}
L.~Wang \textit{et al.},
``Large language model-assisted EDA flow for accelerating circuit design,''
in \textit{Proc. ICCAD}, 2023.

\bibitem{touvron2023}
H.~Touvron \textit{et al.},
``Llama 2: Open foundation and fine-tuned chat models,''
\textit{arXiv preprint}, arXiv:2307.09288, 2023.

\bibitem{brown2020}
T.~B. Brown \textit{et al.},
``Language models are few-shot learners,''
in \textit{Proc. NeurIPS}, 2020, pp.~1877--1901.

\bibitem{hu2021}
E.~J. Hu \textit{et al.},
``LoRA: Low-rank adaptation of large language models,''
in \textit{Proc. ICLR}, 2022.

\end{thebibliography}
\end{document}